\begingroup
\renewcommand{\arraystretch}{1.45}
\fancyhead[L]{\small ZUR THEORIE DER NEUEN TRANSZENDENTEN. III.}
\par \SourceLeaf{462}



\Needspace{27\baselineskip}
und ferner
\[
\begin{gathered}
48.\quad{}
(1-2x^3+2x^{12}- .\,. )+i(1-2x+2x^4- .\,. )\\[8pt]
\quad=1-x\,.\,1-x^3\,.\,1-x^5\mathinner{.\,.\,.} \\[8pt]\quad{}1+i\,.\,1+x^3\,.\,1-x^6\,.\,1+x^9\mathinner{.\,.\,.}\\[8pt]
\qquad\\[8pt]\quad{}\times1-ix\,.\,1-ixx\,.\,\\[8pt]\quad{}1+ix^4\,.\,\\[8pt]\quad{}1+ix^5\,.\,1-ix^7\,.\,1-ix^8\mathinner{.\,.\,.}
\end{gathered}
\]

\Needspace{27\baselineskip}
\[
\begin{gathered}
49.\quad{}
(1+2x^3+2x^{12}+ .\,. )-i(1+2x+2x^4+ .\,. )\\[8pt]
\quad=1+x\,.\,1+x^3\,.\,1+x^5\mathinner{.\,.\,.} \\[8pt]\quad{}1-i\,.\,1-x^3\,.\,1-x^6\,.\,1-x^9\mathinner{.\,.\,.}\\[8pt]
\qquad\\[8pt]\quad{}\times\\[8pt]\quad{}1-ix\,.\,1+ixx\,.\,\\[8pt]\quad{}1-ix^4\,.\,1+ix^5\mathinner{.\,.\,.}
\end{gathered}
\]

\Needspace{27\baselineskip}
\[
\begin{gathered}
50.\quad{}
(1+2x^3+2x^{12}+ .\,. )+i(1+2x+2x^4+ .\,. )\\[8pt]
\quad=1+x\,.\,1+x^3\,.\,1+x^5\mathinner{.\,.\,.} \\[8pt]\quad{}1+i\,.\,1-x^3\,.\,1-x^6\,.\,1-x^9\mathinner{.\,.\,.}\\[8pt]
\qquad\\[8pt]\quad{}\times\\[8pt]\quad{}1+ix\,.\,1-ixx\,.\,\\[8pt]\quad{}1+ix^4\,.\,1-ix^5\mathinner{.\,.\,.}
\end{gathered}
\]


\Needspace{40\baselineskip}
Also aus der Multiplikation von 47 und 48
\[
\begin{gathered}
51.\quad{}
(1-2x^3+2x^{12}-\mathinner{.\,.\,.})^2+(1-2x+2x^4-\mathinner{.\,.\,.})^2\\[8pt]
\quad=2(1-x)^2(1-x^3)^2(1-x^5)^2\mathinner{.\,.\,.}\\[8pt]\quad{}(1+x^3)^2(1-x^6)^2(1+x^9)^2\mathinner{.\,.\,.}\\[8pt]
\qquad\times1+xx\,.\,1+x^4\,.\,1+x^8\,.\,1+x^{10}\mathinner{.\,.\,.}\\[8pt]
\quad=2(1-x)^2(1-x^3)^2(1-x^5)^2\mathinner{.\,.\,.}\\[8pt]\quad{}(1+x^3)^2 1+x^9)^2\mathinner{.\,.\,.}\\[8pt]
\qquad\times(1-x^6)^3(1-x^{12})^2\\[8pt]\quad{}\mathinner{.\,.\,.} 1+xx\,.\,1+x^4\,.\,1+x^6\,.\,1+x^8\mathinner{.\,.\,.}\\[8pt]
\quad=2\frac{[x]^2[x^6]^4}{[xx]^2[x^3]^2[x^{12}]^2}\,.\,\frac{[x^6]^3}{[x^{12}]}\,.\,\frac{[x^4]}{[xx]}
\\[8pt]\quad{}=\frac{2[x]^2[x^4][x^6]^7}{[x^2]^3[x^3]^2[x^{12}]^3}
\end{gathered}
\]


\Needspace{27\baselineskip}
und aus der Multiplikation von 49 und 50
\[
\begin{gathered}
52.\quad{}
(1+2x^3+2x^{12}+ .\,. )^2+(1+2x+2x^4+\mathinner{.\,.\,.})^2\\[8pt]
\quad=2(1+x)^2(1+x^3)^2(1+x^5)^2\mathinner{.\,.\,.}\\[8pt]\quad{}(1-x^3)^2(1-x^6)^2(1-x^9)^2\mathinner{.\,.\,.}\\[8pt]
\qquad\times1+xx\,.\,1+x^4\,.\,1+x^8\,.\,1+x^{10}\mathinner{.\,.\,.}\\[8pt]
\quad=2\frac{[xx]^4}{[x]^2[x^4]^2}[x^3]^2\frac{[x^4][x^6]}{[xx][x^{12}]}
\\[8pt]\quad{}=\frac{2[xx]^3[x^3]^2[x^6]}{[x]^2[x^4][x^{12}]}
\end{gathered}
\]


\Needspace{30\baselineskip}
und der Quotient
\[
\begin{gathered}
53.\quad{}
\frac{(1-2x^3+2x^{12}- .\,. )^2+(1-2x+2x^4- .\,. )^2}
{(1+2x^3+2x^{12}+ .\,. )^2+(1+2x+2x^4+ .\,. )^2}\\[8pt]
\quad=\Bigl(\frac{1-x}{1+x}\Bigr)^2\Bigl(\frac{1-x^3}{1+x^3}\Bigr)^2\Bigl(\frac{1-x^5}{1+x^5}\Bigr)^2\mathinner{.\,.\,.}
\\[8pt]\quad{}\Bigl(\frac{1+x^3}{1-x^3}\Bigr)^2\Bigl(\frac{1+x^9}{1-x^9}\Bigr)^2\mathinner{.\,.\,.}\\[8pt]
\quad=\Bigl(\frac{1-x}{1+x}\Bigr)^2\Bigl(\frac{1-x^5}{1+x^5}\Bigr)^2\Bigl(\frac{1-x^7}{1+x^7}\Bigr)^2\mathinner{.\,.\,.}
\\[8pt]\quad{}=\frac{[x]^4[x^4]^2[x^6]^6}{[x^2]^6[x^3]^4[x^{12}]^2}
\end{gathered}
\]
\SourceLeaf{463}



\Needspace{20\baselineskip}
und das Produkt
\[
\begin{gathered}
54.\quad{}
\{(1-2x^3+ .\,. )^2+(1-2x+ .\,. )^2\}\\[8pt]\quad{}\ \{(1+2x^3+ .\,. )^2-(1+2x+ .\,. )^2\}\\[8pt]
\quad=4\frac{[x^6]^8}{[x^{12}]^4}=4(1-2x^6+2x^{24}- .\,. )^2
\end{gathered}
\]


\Needspace{26\baselineskip}
Aus Formel 23 folgt
\[
\begin{gathered}55.\quad{}
x+x^9+x^{25}+ .\,. \\[8pt]\quad{}=x\,.\,1+x^8\,.\,1+x^{16}\,.\,1+x^{24}\mathinner{.\,.\,.}\\[8pt]\quad{}(1-x^{16}-x^{32}+x^{80}+x^{112}- .\,. )\end{gathered}
\]
\[
\frac32\ \text{Exponent}=\square-1
\]


\Needspace{20\baselineskip}
Aus Formel 26 folgt
\[
\begin{gathered}56.\quad{}
x^3+x^{27}+x^{75}+ .\,. \\[8pt]\quad{}=x\,.\,1+x^8\,.\,1+x^{16}\,.\,1+x^{24}\mathinner{.\,.\,.}\\[8pt]\quad{}(x^2-x^{10}-x^{42}+x^{66}+x^{130}- .\,. )\end{gathered}
\]

\Needspace{15\baselineskip}
\[
\begin{gathered}\text{oder}\quad x^{\frac13}\,.\,1+x^8\,.\,1+x^{16}\,.\,1+x^{24}\mathinner{.\,.\,.}\\[8pt]\quad{}=A,\quad x^2=t^3\quad\text{gesetzt}\end{gathered}
\]

\Needspace{12\baselineskip}
\[
\begin{gathered}55.\quad{}
x+x^9+x^{25}+ .\,. \\[8pt]\quad{}=A(t-t^{25}-t^{49}+t^{121}+t^{169}-\mathinner{.\,.\,.})\end{gathered}
\]

\Needspace{12\baselineskip}
\[
\begin{gathered}56.\quad{}
x^3+x^{27}+x^{75}+ .\,. \\[8pt]\quad{}=A(t^4-t^{16}-t^{64}+t^{100}+t^{196}-\mathinner{.\,.\,.})\end{gathered}
\]


\Needspace{20\baselineskip}
Nun folgt aus der Zerlegung in Faktoren in 24 sehr leicht, wenn man \(x\) durch \(it\)
und \(y\) durch \(i\) ersetzt
\[
\begin{gathered}
it-it^4+it^{16}-it^{25}-it^{49}+ .\,.\\[8pt]
\quad=it\,.\,1-t^3\,.\,1+t^{15}\,.\,1+t^{21}\,.\,1-t^{33}\mathinner{.\,.\,.}\\[8pt]\quad{}1+t^{18}\,.\,1-t^{36}\,.\,1+t^{54}\mathinner{.\,.\,.}
\end{gathered}
\]


\Needspace{23\baselineskip}
Also aus 55---56
\[
\begin{gathered}
57.\quad{}
(x+x^9+x^{25}+ .\,. )-(x^3+x^{27}+x^{75}+ .\,. )\\[8pt]
\quad=x\,.\,1-xx\,.\,1+x^{10}\,.\,1+x^{14}\,.\,1-x^{22}\,.\,1-x^{26}\mathinner{.\,.\,.}\\[8pt]
\qquad\times1+x^{12}\,.\,1-x^{24}\,.\,1+x^{36}\mathinner{.\,.\,.}\\[8pt]\quad{}1+x^8\,.\,1+x^{16}\,.\,1+x^{24}\mathinner{.\,.\,.}
\end{gathered}
\]


\Needspace{23\baselineskip}
Und so sehr leicht
\[
\begin{gathered}
58.\quad{}
(x+x^9+x^{25}+ .\,. )+(x^3+x^{27}+x^{75}+ .\,. )\\[8pt]
\quad=x\,.\,1+xx\,.\,1-x^{10}\,.\,1-x^{14}\,.\,1+x^{22}\,.\,1+x^{26}\mathinner{.\,.\,.}\\[8pt]
\qquad\times1+x^{12}\,.\,1-x^{24}\,.\,1+x^{36}\mathinner{.\,.\,.}\\[8pt]\quad{}1+x^8\,.\,1+x^{16}\,.\,1+x^{24}\mathinner{.\,.\,.}
\end{gathered}
\]

\clearpage
\endgroup
